% 褐矮星（综述）
% license CCBYSA3
% type Wiki

本文根据 CC-BY-SA 协议转载翻译自维基百科\href{https://en.wikipedia.org/wiki/Brown_dwarf}{相关文章}。

\begin{figure}[ht]
\centering
\includegraphics[width=6cm]{./figures/6432b0e71e12ceea.png}
\caption{T 型褐矮星的艺术概念图} \label{fig_Brdw_1}
\end{figure}
\begin{figure}[ht]
\centering
\includegraphics[width=6cm]{./figures/dc30c14f9295e490.png}
\caption{对比：大多数褐矮星的体积仅比木星略大（约大15–20\%）$^{[1]}$，但由于其密度更高，质量仍可达到木星的 80 倍。图像按比例绘制，其中木星的半径约为地球的 11 倍，而太阳的半径约为木星的 10 倍。} \label{fig_Brdw_2}
\end{figure}
褐矮星是亚恒星天体，其质量大于最大的气态巨行星，但小于质量最小的主序星。它们的质量大约为木星的 13–80 倍（($M_J$)）$^{[2]}$——不足以在其核心中维持将氢聚变为氦的核反应，但又足够巨大，可以通过氘（$^2\mathrm{H}$）的聚变释放一定的光和热。氘是氢的一种同位素，含有一个中子和一个质子，其聚变所需温度更低。质量最大的褐矮星（$(>65,M_J)$）还能够进行锂（$(^7\mathrm{Li})$）的聚变。

天文学家按照光谱型对自发光天体进行分类，这一分类与表面温度密切相关。褐矮星分布在 M 型（2100–3500 K）、L 型（1300–2100 K）、T 型（600–1300 K）和 Y 型（<600 K） 光谱类别中 $^{[3][4]}$。由于褐矮星不会发生稳定的氢聚变，它们会随着时间推移逐渐冷却，并在演化过程中依次进入更晚的光谱型。

“褐矮星”中的“褐色”原意是指介于红色与黑色之间的一种颜色 $^{[5]}$。以肉眼观察，大多数褐矮星呈现为品红色，而其他褐矮星则会因温度不同而显示出不同的颜色 $^{[3][6]}$。褐矮星可能是完全对流的，内部不存在明显的分层结构或随深度变化的化学分异 $^{[7]}$。

尽管早在 20 世纪 60 年代就有人从理论上预言了褐矮星的存在，但直到 1994 年，第一批明确无误的褐矮星才被发现 $^{[8]}$。由于褐矮星的表面温度相对较低，它们在可见光波段并不明亮，而主要在红外波段辐射能量。随着红外探测设备能力的提升，天文学家已经识别出数以千计的褐矮星。目前已知距离最近的褐矮星位于 Luhman 16 系统，这是一个由 L 型和 T 型褐矮星组成的双星系统，距离太阳约 6.5 光年（2.0 秒差距）。在距离上，Luhman 16 是继半人马座 $\alpha$ 星系和巴纳德星之后，距离太阳第三近的恒星系统。
\subsection{历史}
\begin{figure}[ht]
\centering
\includegraphics[width=6cm]{./figures/16619a30e2a82cab.png}
\caption{较小的天体是葛利泽 229B（Gliese 229B），其质量约为木星的 20 至 50 倍，绕恒星葛利泽 229 运行。它位于天兔座，距离地球约 19 光年。} \label{fig_Brdw_3}
\end{figure}
\subsubsection{早期理论研究}
\begin{figure}[ht]
\centering
\includegraphics[width=6cm]{./figures/a5a243a9e6a937ce.png}
\caption{行星、褐矮星、恒星（非按比例）} \label{fig_Brdw_4}
\end{figure}
20 世纪 60 年代，Shiv Kumar 提出了后来被称为褐矮星的天体的存在理论；它们最初被称为黑矮星$^{[9]}$，这一名称用于指在太空中自由漂浮、质量不足以维持氢聚变的暗弱亚恒星天体。然而：\\

(a) “黑矮星”一词此前已被用来指冷却后的白矮星；\\
(b) 红矮星能够进行氢聚变；\\
(c) 这些天体在其生命早期可能在可见波段仍然具有一定亮度。\\

因此，人们提出了其他名称来称呼这些天体，包括 planetar 和 substar。1975 年，Jill Tarter 在加州大学伯克利分校的博士论文中，首次建议使用“褐矮星（brown dwarf）”这一术语，用“褐色”来表示一种“介于红色与黑色之间”的颜色，意指这些矮星看起来昏暗、黯淡、无光泽，尽管它们并非真正呈褐色$^{[5][10][11]}$。

“黑矮星”这一术语至今仍用于指已经冷却到几乎不再发出可探测光辐射的白矮星。然而，即便是最低质量的白矮星，要冷却到这一温度所需的时间也被计算为超过宇宙当前的年龄，因此这类天体预计至今尚未存在$^{[12]}$。

关于最低质量恒星及氢燃烧下限的早期理论认为：第一星族天体若质量小于 $0.07,M_\odot$，或第二星族天体若质量小于 $0.09,M_\odot$，将永远不会经历正常的恒星演化过程，而会成为完全简并的恒星$^{[13]}$。由此产生的褐矮星有时被称为“失败的恒星”$^{[14]}$。首个自洽的氢燃烧最小质量计算结果证实，第一星族天体的下限约为 $0.07-0.08,M_\odot$,$^{[15][16]}$。
\subsubsection{氘聚变}
20 世纪 80 年代末，人们发现质量低至 $0.013,M_\odot$（约 $13.6,M_J$）的天体也可以发生氘聚变，同时褐矮星冷却外层大气中尘埃形成的影响也被认识到，这些发现对上述理论提出了挑战。然而，这类天体极难被发现，因为它们几乎不在可见光波段发射辐射，其最强的辐射位于红外（IR）波段。当时的地基红外探测器精度不足，难以可靠识别褐矮星。

此后，研究者采用多种方法开展了大量搜索工作，包括：对场星的多色成像巡天；对主序矮星和白矮星暗弱伴星的成像巡天；对年轻星团的调查；以及通过径向速度监测寻找近距离伴星。
\subsubsection{GD 165B 与 L 型}

多年来，发现褐矮星的尝试一直未获成功。直到 1988 年，在一次针对白矮星的红外巡天中，人们发现白矮星 GD 165 周围存在一个暗弱伴星。该伴星 GD 165B 的光谱呈现出非常红且难以解释的特征，既不符合低质量红矮星的典型特征，也不同于当时已知的最冷 M 型矮星。这表明，GD 165B 必须被归类为一种比已知 M 型矮星更冷的天体。

在接近十年的时间里，GD 165B 一直是独一无二的例子。直到 1997 年 2MASS（Two Micron All-Sky Survey）启动后，人们才发现了大量在颜色和光谱特征上与其相似的天体。

如今，GD 165B 被认为是后来被称为 L 型矮星的一类天体的原型$^{[17][18]}$。

尽管当时发现如此冷的矮星具有重大意义，但学界一度争论 GD 165B 究竟应被归类为褐矮星，还是仅仅是一颗极低质量恒星，因为在观测上二者极难区分$^{[19][20]}$。

在 GD 165B 被发现之后不久，又陆续有其他褐矮星候选体被报告。然而，其中大多数未能通过检验，因为它们的大气中缺乏**锂元素**，表明它们实际上是恒星。真正的恒星会在约 $10^8$ 年左右燃尽其锂，而褐矮星（尽管其温度和光度可能与恒星相近）则不会发生这一过程。因此，在年龄超过 $100,\mathrm{Myr}$ 的天体大气中探测到锂元素，几乎可以确定该天体是一颗褐矮星。
\subsubsection{Teide 1 与 M }
首个被确认的 **M 型褐矮星** 于 1994 年由西班牙天体物理学家 Rafael Rebolo（课题组负责人）、María Rosa Zapatero-Osorio 和 Eduardo L. Martín 发现$^{[21]}$。该天体位于昴宿星团这一疏散星团中，被命名为 **Teide 1**。相关发现论文于 1995 年 5 月投稿至《自然》，并于 1995 年 9 月 14 日正式发表$^{[22][23]}$。在该期《自然》杂志的封面上，醒目标题写道：“**褐矮星被发现，官方确认**”。

Teide 1 是 IAC 团队于 1994 年 1 月 6 日在特内里费岛泰德天文台，使用 80 厘米口径望远镜（IAC 80）获取的图像中发现的；其光谱则于 1994 年 12 月首次由拉帕尔马岛罗克·德·洛斯·穆查乔斯天文台的 4.2 米 **威廉·赫歇尔望远镜** 观测获得。由于 Teide 1 属于年轻的昴宿星团，其距离、化学组成和年龄得以较为准确地确定。利用当时最先进的恒星与亚恒星演化模型，研究团队估计 Teide 1 的质量为 (55 \pm 15,M_J)$^{[24]}$，低于恒星质量下限。此后，Teide 1 成为年轻褐矮星研究中的一个重要参照对象。

从理论上讲，质量低于$65,M_J) 的褐矮星在其整个演化过程中都无法通过热核聚变燃烧锂。这一事实构成了锂检验原理的基础之一，用以判断低光度、低表面温度天体是否具有亚恒星性质。

1995 年 11 月，使用 Keck I 望远镜获得的高质量光谱数据显示，Teide 1 仍然保留着形成昴宿星团恒星的原始分子云中的初始锂丰度，从而证明其核心中不存在热核聚变过程。这些观测结果不仅完全确认了 Teide 1 是一颗褐矮星，也验证了光谱锂检验方法的有效性$^{[25]}$。

在相当一段时间内，Teide 1 是通过直接观测所识别的、太阳系外已知质量最小的天体。此后，人们已鉴认出超过 1,800 颗褐矮星$^{[26]}$，其中一些距离地球非常近，例如 Epsilon Indi Ba 与 Bb——一对与一颗类太阳恒星引力束缚、距离太阳约 12 光年的褐矮星系统$^{[27]}$，以及 Luhman 16——一个距离太阳仅 6.5 光年的褐矮星双星系统。
\subsubsection{Gliese 229B 与 T 型}
首个 T 型褐矮星于 1994 年由加州理工学院的天文学家 Shrinivas Kulkarni、中岛敬史（Tadashi Nakajima）、Keith Matthews 和 Rebecca Oppenheimer，以及约翰斯·霍普金斯大学的科学家 Samuel T. Durrance 和 David Golimowski 共同发现$^{[28]}$。该天体于 1995 年被确认是 Gliese 229 的一颗亚恒星伴星，即 Gliese 229b。Gliese 229b 与 Teide 1 一同构成了最早两例清晰的褐矮星观测证据。二者均在 1995 年通过探测到 670.8 nm 的锂谱线得到确认，其中 Gliese 229b 的温度和光度显著低于恒星范围。

其近红外光谱清楚地显示出 2 微米处的甲烷吸收带，这一特征此前仅在巨行星大气以及土星卫星泰坦的大气中被观测到。甲烷吸收并不可能出现在任何主序恒星的温度条件下。正是这一发现促成了比 L 型矮星更冷的新光谱类型——T 型矮星的确立，而 Gliese 229b 则被视为该类型的原型。
\subsection{理论}
\begin{figure}[ht]
\centering
\includegraphics[width=6cm]{./figures/851acc2df0619164.png}
\caption{} \label{fig_Brdw_5}
\end{figure}